% Additive proof for the original 15671-variable polynomial. No circuit changes.
\subsection{The seven-witness parity predicate and its finite correction}
\label{subsec:original-seven-parity}
Put $B=17b$. The predicate used here is exactly the existence of
$t,U,e,p,q,s,h\in\mathbb Q$ satisfying
\begin{align}
 U(1+bt^2)&=b(1-bt^2),\label{eq:seven-param}\\
 p^2&=B(1-e)(1+3e),\label{eq:seven-first}\\
 h&=e(B-U^2),\label{eq:seven-h}\\
 Bq^2&=(s^2-Bh)(s^2+3Bh).\label{eq:seven-second}
\end{align}
Denote it by $\Psi(b)$. No nonzero conditions are added to this formula.

\begin{lemma}[The exact one-sided assertion]
\label{lem:original-seven-parity-soundness}
If $b\ne0$ and $\Psi(b)$, then $v_\ell(b)$ is not a negative odd
integer for any prime $\ell\notin\{2,3,17\}$.
\end{lemma}
\begin{proof}
Suppose $k=v_\ell(b)<0$ is odd. Equation~\eqref{eq:seven-param}
forces $1+bt^2\ne0$: if it vanished, its right side would be $2b\ne0$.
If $t=0$, then $U=b$. Otherwise $v_\ell(bt^2)$ is odd and nonzero,
so $(1-bt^2)/(1+bt^2)$ is a unit. Thus $v_\ell(U)=k$ in all cases,
and $v_\ell(B-U^2)=2k$.

If $e=0$, equation~\eqref{eq:seven-first} gives $p^2=B$, impossible.
Write $j=v_\ell(e)$. If $j<0$, the right side of
\eqref{eq:seven-first} has valuation $k+2j$, which is odd.
If $j>0$, it has valuation $k$, again odd. Consequently $j=0$.
This includes the possibilities $p=0$, $e=1$, and $e=-1/3$.
Now $h\ne0$ and $v_\ell(Bh)=3k$ is odd. If $s\ne0$, its square
has even valuation, different from $3k$, so the two factors on the
right of~\eqref{eq:seven-second} have the same finite valuation.
Their product has even valuation and is nonzero. The same conclusion
holds when $s=0$, since the product is $-3(Bh)^2$.
The left side is either zero or has odd valuation, a contradiction.
\end{proof}

The omitted guards therefore cause no loss of the asserted soundness.
At $b=0$, the equations require $U=p=h=s=0$ and permit arbitrary
$t,e,q$. For $b\ne0$, the six-witness source formula is recovered
exactly by $e_{\rm old}=Be$, $p_{\rm old}=Bp$, $q_{\rm old}=Bq$,
with $s,t,U$ unchanged. Its first curve residual is $B^2$ times
\eqref{eq:seven-first}'s residual and its second is $B$ times
\eqref{eq:seven-second}'s residual after~\eqref{eq:seven-h}.
Division by $B$ reverses this transformation.

\begin{lemma}[Simultaneous realization from the existing restricted converse]
\label{lem:original-seven-simultaneous-realization}
Assume the manuscript's restricted converse $\mathsf{RC}$ for
$E_l:y^2=x(x-l)(x+3l)$ with $l>0$ squarefree and $l\equiv11\pmod{12}$.
Let $\beta,\theta>0$ be squarefree, with disjoint supports avoiding
$2,3$, such that $\beta\equiv17\pmod{24}$,
$\theta\equiv1\pmod{24}$, and each is a local square at every prime
dividing the other. Then there is a positive squarefree $n$, prime
to $6\theta$, for which both $E_n$ and $E_{n\theta}$ have rational
points with nonzero abscissa in $\beta\mathbb Q^{\times2}$.
The case $\theta=1$ is allowed.
\end{lemma}
\begin{proof}
Here are the descent steps, so that no separate realization theorem
is being assumed. Write $B_0=\operatorname{supp}(\beta)$,
$J=\operatorname{supp}(\theta)$, and $d=|B_0|\ge1$.
For a support prime $q_i$, put $s_i=[-1/q_i]$, $\lambda_i=[3/q_i]$,
and define over $\mathbb F_2$
\[
 (Kv)_i=\sum_{j\ne i}[q_j/q_i](v_j+v_i).
\]
Add $d-1$ prime slots $I$ of type $(s,\lambda)=(0,0)$ and one slot
$q_0$ of type $(1,1)$. Let the $B_0$-by-$I$ symbol matrix have columns
$e_i+e_d$ for $1\le i<d$. Set all other symbols incident with $I$
to zero, prescribe $[q_0/q_b]$ as the corresponding row sum, and
obtain reverse symbols by reciprocity. The $q_0$--$J$ symbols remain
free. Mutual squareness gives zero row sums for both directed
$B_0$--$J$ blocks; hence $K\mathbf1_{B_0}=0$ on both supports.

Choose $a$ with ones at one fixed member of $B_0$ and at $q_0$.
The matrix $K^\sharp=K+\lambda a^t$ kills $\mathbf1$ on both sides.
After deleting $q_0$, its smaller-support matrix has the form
\[
 M_0=\begin{pmatrix}D&C\\C^t&0\end{pmatrix},
 \qquad \mathbf1^tD\mathbf1=1.
\]
The columns of $C$ span the even-sum hyperplane. If $M_0(x,y)^t=0$,
then $x=c\mathbf1$; multiplication of the first equation by
$\mathbf1^t$ gives $c=0$, and injectivity of $C$ gives $y=0$.
This also proves the case $d=1$, where $M_0=(1)$.
Adding $J$ preserves $M_0$. Each free $q_0$--$J$ bit toggles just
one new diagonal entry. The enlarged determinant is multilinear in
these bits, with full-product coefficient $\det M_0=1$, so some
Boolean choice makes it nonzero. CRT and Dirichlet's theorem realize
the prescribed symbols by distinct fresh primes, with types imposed
modulo $12$. Thus $n=\prod_{B_0\cup I\cup\{q_0\}}q$ and $n\theta$
are squarefree and congruent to $11$ modulo $12$.

The projection $T=1-\mathbf1_{B_0}a^t$ satisfies $K=K^\sharp T$.
Nonsingularity of the deleted matrix gives
$\ker K^\sharp=\langle\mathbf1\rangle$, and therefore
\[
 \ker K=\langle\mathbf1,\mathbf1_{B_0}\rangle,
 \qquad \lambda\notin\operatorname{im}K
\]
on both supports. For the latter assertion, $Kx=\lambda$ would
make $Tx-\mathbf1_{B_0}$ constant, contradicted by applying $a^t$.

For clarity, the local descent needed here is the manuscript's full
two-descent in coordinates $(x+3l,x-l)$. The two torsion classes are
$(l,1)$ and $(3,-l)$. At an odd good prime, the local image consists
of the four pairs with even valuations; at a prime dividing $l$,
the two torsion classes span it by their independent valuation
vectors. At $3$ it is the set of all first-coordinate classes with
square second coordinate. At $2$ both coordinate valuations are
even: the only problematic case $x-l=2u$, $u$ odd, would give
$(y/2)^2\equiv l\pmod8$, impossible for $l\equiv3\pmod4$.
When both differences are divisible by $4$, their difference $4l$
and the cubic equation give even valuations for both. At infinity
the first coordinate is positive, and the positive-second-coordinate
subgroup has index two. Its possible representatives consequently are
$(3^\delta\prod q_i^{u_i},\prod q_i^{v_i})$ with necessary equations
\[
 Ku+\operatorname{diag}(\lambda)v=\delta\lambda,
 \quad (K+\operatorname{diag}(s))v=0,
 \quad (\lambda+s)^tv=0.
\]
Reciprocity gives $K+K^t=ss^t+\operatorname{diag}(s)$.
Hence $\ker(K+\operatorname{diag}(s))=\ker K^t\cap s^\perp$
is a line. The functional $\lambda^t$ is nonzero on this line,
because $\lambda\notin\operatorname{im}K$ and
$\lambda^t\mathbf1=0$, whereas $s^t\mathbf1=1$.
The displayed equations force $v=0$, $\delta=0$, and $u\in\ker K$.
There are at most eight Selmer classes.
The class $(\beta,1)$ is actually local everywhere: it is the identity
at $2$ and infinity, the first torsion class at $3$, satisfies the
torsion comparison at support primes by $K\mathbf1_{B_0}=0$, and
is unramified at the other odd primes. It is independent of the two
torsion classes since $B_0$ is nonempty and $q_0\notin B_0$.
Thus each two-Selmer group has dimension exactly three.

The existing dimension-three lemma gives full two-primary Selmer
corank one. Explicitly, the Kummer sequence gives
$r+\dim\Sha[2]=1$; writing
$\Sha[2^\infty]=(\mathbb Q_2/\mathbb Z_2)^t\oplus F$, the
alternating Cassels--Tate pairing on the finite quotient makes
$\dim F[2]$ even, so $F=0$ and $r+t=1$.
Apply $\mathsf{RC}$ to $n$ and $n\theta$, and the manuscript's
published forward analytic-rank-one theorem. Both curves have rank
one and finite $\Sha$, hence $\Sha[2]=0$. Kummer exactness now
represents $(\beta,1)$ by actual rational points. These are not
two-torsion, so their abscissas are nonzero of class $\beta$.
When $\theta=1$ the same point can be used twice.
\end{proof}

\begin{lemma}[The completeness scope actually used]
\label{lem:original-seven-parity-completeness}
Assume only the existing $\mathsf{RC}$ and the published descent and
forward-rank facts used above. If $b>0$ has square reduced denominator
and is a square in $\mathbb Q_2,\mathbb Q_3,\mathbb Q_{17}$, then
$\Psi(b)$ holds.
\end{lemma}
\begin{proof}
Choose $t\in\mathbb Q$ close to a solution of $bt^2=1$ at infinity
and at $2,3,17$. In~\eqref{eq:seven-param} this makes $U$ arbitrarily
small at those places. We can and do require
$V=B(B-U^2)>0$ and $V$ locally square at all three primes.
Explicit sufficient conditions are $U^2<B$ in $\mathbb R$,
$v_2(U^2/B)\ge3$, and $v_3(U^2/B),v_{17}(U^2/B)\ge1$.
The denominator tends locally to $2\ne0$ and is positive globally.

Let $\beta,\theta$ be the positive squarefree representatives of
$B,V$. The number $B=17b$ is square at $2$, a nonsquare unit class
at $3$, and has odd valuation at $17$. Thus
$\beta\equiv17\pmod{24}$ and $\beta>1$;
$\theta\equiv1\pmod{24}$ and $17\nmid\theta$.
At any prime $\ell\mid\beta$ outside $2,3,17$, the valuation of
$b$ is odd and therefore positive, since its reduced denominator is
square. The same quotient calculation as in soundness gives
$v_\ell(U)=v_\ell(b)$; hence $1-U^2/B$ is a square unit.
This proves that $V$, and hence $\theta$, is square at every prime
of $\beta$, including $17$ by construction. The supports are disjoint.
Conversely, if $\ell\mid\theta$, then $\ell\notin\{2,3,17\}$
and $v_\ell(B)$ is even. Odd valuation of $B(B-U^2)$ forces
$v_\ell(B)=2v_\ell(U)$ and strict cancellation, so
$B/U^2\in1+\ell\mathbb Z_\ell$ is square. This proves the other
mutual-square condition. If $U=0$, then $\theta=1$ and this last
argument is unnecessary.

Apply the preceding lemma. Write $V=\theta c^2$ and rescale the
second point by $(x,y)\mapsto(c^2x,c^3y)$ to put it on $E_{nV}$.
If the first point has abscissa $Ba^2$, divide both curve parameters
by $a^2$ and the point coordinates by $a^2,a^3$. The resulting
points have abscissas $B,Bs^2$ with $s\ne0$, on $E_{e_0}$ and
$E_{e_0V}$ respectively. Write their ordinates as $y_1,y_2$.
Set
\[
 e=e_0/B,\qquad p=y_1/B,\qquad
 q=y_2/(B^2s),\qquad h=e(B-U^2).
\]
Direct substitution verifies all four equations. Every division
made in choosing these witnesses is by a nonzero number; none is
an additional equation or guard in the original polynomial.
\end{proof}

\paragraph{Finite correction and the unguarded norm root.}
At $2$ use valuation parity followed by unit coordinates $[-1],[5]$;
at $3$ and $17$ use valuation parity followed by the nonsquare-unit bit.
The columns for $2,3,5,7,11,13,17$ are exactly
\[
\begin{pmatrix}
1&0&0&0&0&0&0\\0&1&0&1&1&0&0\\0&1&1&0&1&1&0\\
0&1&0&0&0&0&0\\1&0&1&0&1&0&1\\0&0&0&0&0&0&1\\
0&1&1&1&1&0&0
\end{pmatrix}.
\]
The accompanying exact computation verifies rank seven and all 128
distinct subset signatures. Rational solutions of $e^2-e=0$ are
exactly $0,1$. Therefore the original eight bits and two products
\[
 A=(1+e_2)(1+2e_3)(1+12e_{13})(1+16e_{17}),\quad
 D=(1+4e_5)(1+6e_7)(1+10e_{11})(1-2\eta)
\]
give exactly $AD=(-1)^\eta\prod p^{e_p}$, always nonzero.
For nonzero integral $r$, choose $\delta=\operatorname{sign}(r)$
and the unique subset product $\mu$ having the same local signature
as $\delta r$. Taking $AD=\delta\mu$ makes $ADr$ positive and
locally square at $2,3,17$.

Choose $J_0$ sufficiently deep at all places over primes dividing
$\mu r$, then $S$ in a sufficiently divisible subgroup satisfying
the anchor. Since $f=z^2$ times a formal unit, the coefficients of
$f(J_0),f(S),z(S)/z(J_0)$ can attain any prescribed finite lower
valuation bounds. This remains true at ramified and bad primes by
working in sufficiently deep formal subgroups. Thus at $\ell\mid r$
all normalized leaves are integral, and at $\ell\mid\mu$ with
$\ell\nmid r$ all nonconstant leaves can be made divisible by $\ell$.
The norm tree $N_r(a,b)=a^2+rb^2$ has its constant leaf $1$ on
the all-left path; its root $W$ is a unit at these primes.
In the actual application $r=\beta(\beta-1)(\beta-2401)
(32\beta+343)(2\beta-1)$ with $\beta=X^2+2$ and $X\in\mathbb Z$.
It is nonzero and $|r|>1$ since $X^2\ne2399$ and $\beta\ge2$.
Consequently $W\ne0$, and $b=ADrW^2$ is positive, locally square
at $2,3,17$, and has square reduced denominator. Completeness applies.

For soundness no global guard $W\ne0$ is needed. At any prime where
$v_\ell(r)$ is odd, the two nonzero summands of $a^2+rb^2$ have
different valuation parity. Such a norm vanishes only if both inputs
vanish. Recursing to the constant leaf $1$ proves $W\ne0$ at every
prime relevant to an odd-contact argument. Outside the fixed
exceptional set, $AD$ is a unit and
$v_\ell(ADrW^2)=v_\ell(r)+2v_\ell(W)$.
The exact original norm argument therefore applies unchanged. If
$r$ has no relevant odd valuation, there is no local conclusion to
prove. No guard or witness has been inserted, and the original
polynomial and coefficient bound are unchanged.
